% ============================================================
% 65_code_style.tex -- Informatik-Code-Layout
% ============================================================

\usepackage{listings}
\usepackage{inconsolata}
\tcbuselibrary{listings}

\lstset{
  basicstyle=\ttfamily,
  columns=flexible,
  breaklines=true,
  breakatwhitespace=true,
  tabsize=2,
  showstringspaces=false,
}

\lstdefinestyle{CodeExpert}{
  language=Python,
  basicstyle=\ttfamily\linespread{0.82}\selectfont\color{ce_white},
  numbers=none,
  frame=none,
  aboveskip=\ZSFcodeMiniSkip,
  belowskip=\ZSFcodeMiniSkip,
  xleftmargin=0pt,
  xrightmargin=0pt,
  keywordstyle=\bfseries\color{ce_pink},
  commentstyle=\color{ce_green},
  stringstyle=\color{ce_yellow},
  identifierstyle=\color{ce_white},
  breaklines=true,
  breakatwhitespace=true,
  literate={->}{{$\rightarrow$}}2 {ü}{{\"u}}1 {ä}{{\"a}}1 {ö}{{\"o}}1 {Ü}{{\"U}}1 {Ä}{{\"A}}1 {Ö}{{\"O}}1 {—}{{---}}1 {−}{{-}}1 {„}{{\glqq}}1 {“}{{\grqq}}1 {”}{{''}}1 {’}{{'}}1 {–}{{--}}1 {→}{{$\rightarrow$}}1
}

\tcbset{zsfcodeboxbase/.style={
  enhanced jigsaw,
  colback=code_bg,
  colframe=code_rule,
  boxrule=\ZSFcodeFrameRule,
  boxsep=0pt,
  arc=\ZSFcornerBox, outer arc=\ZSFcornerBox,
  left=\ZSFcodeInnerXPad, right=\ZSFcodeInnerXPad,
  top=\ZSFcodeInnerYPad, bottom=\ZSFcodeInnerYPad,
  before skip=\ZSFboxBeforeSkip,
  after skip=\ZSFboxAfterSkip,
  before upper={\parskip=0pt\relax},
}}
\tcbset{zsfcodeboxtitled/.style={
  zsfcodeboxbase,
  colbacktitle=code_title_bg,
  coltitle=code_title_text,
  titlerule=0pt,
  zsftitlebar,
}}

\NewTColorBox{codebox}{o O{}}{
  zsfcodeboxbase,
  IfValueT={#1}{%
    zsfcodeboxtitled,
    title={\strut #1\if\relax\detokenize{#2}\relax\else\ZSFTitleTag{#2}\fi}},
}

\tcbset{zsf@codebox@split@base/.style={
  enhanced jigsaw,
  colback=code_bg,
  colframe=code_rule,
  leftrule=\ZSFcodeFrameRule,
  rightrule=\ZSFcodeFrameRule,
  toprule=0pt,
  bottomrule=0pt,
  arc=\ZSFcornerBox, outer arc=\ZSFcornerBox,
  boxsep=0pt,
  left=\ZSFcodeInnerXPad, right=\ZSFcodeInnerXPad,
  top=\ZSFcodeInnerYPad, bottom=\ZSFcodeInnerYPad,
  before skip=0pt,
  after skip=\ZSFspaceXXS,
  before upper={\parskip=0pt\relax},
}}
\tcbset{zsf@codebox@split@titled/.style={
  zsf@codebox@split@base,
  colbacktitle=code_title_bg,
  coltitle=code_title_text,
  titlerule=0pt,
  zsftitlebar,
}}
\NewTColorBox{codeboxfirst}{o}{
  zsf@codebox@split@base,
  IfValueT={#1}{zsf@codebox@split@titled,title={\strut #1}},
  toprule=\ZSFcodeFrameRule,
  sharp corners=south,
  before skip=\ZSFboxBeforeSkip,
}

\NewTColorBox{codeboxmid}{o}{
  zsf@codebox@split@base,
  IfValueT={#1}{zsf@codebox@split@titled,title={\strut #1}},
}

\NewTColorBox{codeboxlast}{o}{
  zsf@codebox@split@base,
  IfValueT={#1}{zsf@codebox@split@titled,title={\strut #1}},
  bottomrule=\ZSFcodeFrameRule,
  sharp corners=north,
  after skip=\ZSFboxAfterSkip,
}

% Kompakter Katalog für parametrisierbare Schleifenmuster.
% Eine Box bündelt mehrere Muster derselben Kategorie; jedes endet mit Theta.
\newif\ifzsfloopcatalogfirst
\NewTColorBox{loopcatalog}{}{
  zsfcodeboxbase,
  left=\ZSFspaceS,
  right=\ZSFspaceS,
  top=0pt,
  bottom=\ZSFspaceXS,
  before upper={\parskip=0pt\relax\zsfloopcatalogfirsttrue},
}

\newcommand{\LoopPatternTitle}[1]{%
  \par
  \ifzsfloopcatalogfirst
    \zsfloopcatalogfirstfalse
  \else
    {\color{code_rule}\hrule height\ZSFborderThin}\addvspace{\ZSFspaceXS}%
  \fi
  \noindent{\ZSFfontBoxTitle #1}\par\nobreak}
\newcommand{\LoopPatternResult}[1]{%
  \nointerlineskip
  \hbox to \linewidth{\hfil\smash{\raisebox{\dimexpr\baselineskip-\ZSFcodeMiniSkip\relax}{\eqbox{\Theta\!\left(#1\right)}}}}\par}
